\chapter{Errores y sesgos: lo que el registro no cuenta}\label{cap:errores}

\objetivos{
\begin{itemize}
  \item Distinguir errores de omisión y de comisión y calcular sus tasas.
  \item Diferenciar un error aleatorio de un sesgo sistemático.
  \item Enumerar las fuentes de omisión, de falsas alarmas y de conteo múltiple de cada sensor.
  \item Reconocer las fuentes térmicas persistentes que no son fuego de vegetación.
  \item Decidir qué errores afectan a la armonización y cuáles no.
\end{itemize}}

\section{Ningún detector es perfecto}

Los capítulos anteriores mostraron cómo se llega de las llamas a un punto en un archivo. Cada
paso puede fallar en dos sentidos: dejar pasar un fuego real o inventar uno que no existe. Este
capítulo ordena esos fallos, porque la armonización consiste, en el fondo, en entender en qué se
diferencian los errores de los dos sensores.

\section{Omisión y comisión}

Comparemos lo que dice el sensor con la realidad. Para cada píxel hay cuatro posibilidades,
recogidas en la \termino{matriz de confusión} de la tabla~\ref{tab:confusion}.

\begin{table}[ht]
\centering
\renewcommand{\arraystretch}{1.4}
\begin{tabular}{l|cc}
 & \textbf{Hay fuego} & \textbf{No hay fuego} \\
\hline
\textbf{El sensor detecta} & acierto (VP) & \cellcolor{fuegoClaro}falsa alarma (FP) \\
\textbf{El sensor no detecta} & \cellcolor{brasaClaro}fuego omitido (FN) & rechazo correcto (VN) \\
\end{tabular}
\caption{Matriz de confusión de un detector de fuego. VP: verdadero positivo; FP: falso
positivo; FN: falso negativo; VN: verdadero negativo.}
\label{tab:confusion}
\end{table}

\begin{definicion}{error de omisión y error de comisión}
El \textbf{error de omisión} es la fracción de los fuegos reales que el sensor no detectó:
\[
  \text{omisión} = \frac{FN}{VP + FN}.
\]
El \textbf{error de comisión} es la fracción de las detecciones que no corresponden a un fuego
real:
\[
  \text{comisión} = \frac{FP}{VP + FP}.
\]
\end{definicion}

Son cosas distintas y se miden respecto de totales distintos: la omisión, respecto de lo que
había; la comisión, respecto de lo que se detectó. El complemento de la omisión se llama
\emph{probabilidad de detección} o sensibilidad; el de la comisión, \emph{precisión}.

\begin{ejemplo}{tasas de error}
En una zona de validación había 200 fuegos confirmados con imágenes de alta resolución. Un
sensor reportó 150 detecciones, de las cuales 144 coincidían con fuegos reales. Calcula sus
errores.
\solucion $VP = 144$, $FP = 150 - 144 = 6$, $FN = 200 - 144 = 56$.
\[
 \text{omisión} = \frac{56}{200} = 28\,{}\%, \qquad \text{comisión} = \frac{6}{150} = 4\,{}\%.
\]
El sensor casi nunca se equivoca cuando dice que hay fuego, pero se le escapa más de la cuarta
parte de los fuegos.
\end{ejemplo}

El ejemplo refleja la situación real. Los algoritmos de fuego están diseñados para ser
\textbf{conservadores}: se prefiere perder fuegos pequeños a inundar los mapas de falsas
alarmas. La comisión de MODIS y de VIIRS es de unos pocos puntos porcentuales o menos; la
omisión, en cambio, es grande, muy variable y depende sobre todo del tamaño de los fuegos.

\begin{intuicion}[La omisión no es un número, es una curva]
No tiene sentido decir <<MODIS omite el 30\,{}\% de los fuegos>>. Omite casi todos los de
50~m$^2$ y casi ninguno de 10~hectáreas. Lo que existe es una \emph{curva de probabilidad de
detección} en función del tamaño (figura~\ref{fig:pdet}), y la omisión total en una región
depende de qué tamaños de fuego abundan allí.
\end{intuicion}

\begin{figure}[ht]
\centering
\begin{tikzpicture}
\begin{axis}[width=13cm,height=6.8cm,xmode=log,xmin=5,xmax=100000,ymin=0,ymax=1.08,
  xlabel={Área del fuego (m$^2$)},ylabel={Probabilidad de detección},
  ytick={0,0.5,1},yticklabels={0,{0,5},1},
  legend style={at={(0.97,0.06)},anchor=south east,font=\small,draw=ceniza!40},
  grid=major,grid style={ceniza!15},samples=120,domain=5:100000]
\addplot[viirs,very thick] {1/(1+exp(-3.2*(ln(x)/ln(10)-ln(40)/ln(10))))};
\addplot[modis,very thick] {1/(1+exp(-3.2*(ln(x)/ln(10)-ln(300)/ln(10))))};
\addplot[modis,very thick,dashed] {1/(1+exp(-3.2*(ln(x)/ln(10)-ln(2900)/ln(10))))};
\legend{VIIRS 375 m (nadir),MODIS (nadir),MODIS (borde de la franja)}
\fill[brasa,opacity=0.15] (axis cs:40,0) rectangle (axis cs:300,1.08);
\node[font=\footnotesize,align=center,fill=white,inner sep=2pt] at (axis cs:100,0.93) {solo\\ VIIRS};
\end{axis}
\end{tikzpicture}
\caption{Forma esperada de la probabilidad de detección según el tamaño del fuego (curvas
esquemáticas basadas en la proporcionalidad de la ecuación~\eqref{eq:amin}; no son mediciones).
Los fuegos de la banda sombreada los detecta VIIRS y no MODIS.}
\label{fig:pdet}
\end{figure}

\section{Error aleatorio y sesgo}

Hay otra distinción, igual de importante, que atraviesa todo el capítulo.

\begin{definicion}{error aleatorio y sesgo}
Un \textbf{error aleatorio} varía de una medición a otra sin dirección preferente; al promediar
muchas mediciones tiende a cancelarse. Un \textbf{sesgo}, o error sistemático, empuja siempre en
el mismo sentido; promediar no lo reduce.
\end{definicion}

\begin{figure}[ht]
\centering
\begin{tikzpicture}
  \foreach \x/\t/\cx/\cy/\s in {0/{Poco sesgo, poca dispersión}/0/0/0.18,4.6/{Poco sesgo, mucha dispersión}/0/0/0.75,9.2/{Mucho sesgo, poca dispersión}/0.75/0.6/0.18}{
    \begin{scope}[shift={(\x,0)}]
      \foreach \r in {1.6,1.1,0.6}{\draw[ceniza!70] (0,0) circle (\r);}
      \fill[ceniza] (0,0) circle (0.06);
      \pgfmathsetseed{7}
      \foreach \i in {1,...,12}{\fill[fuego] ({\cx+\s*rand},{\cy+\s*rand}) circle (0.07);}
      \node[font=\footnotesize,align=center,text width=4cm] at (0,-2.15) {\t};
    \end{scope}
  }
\end{tikzpicture}
\caption{Dispersión y sesgo. El centro de la diana es el valor verdadero. Promediar muchos
tiros arregla el segundo caso, pero no el tercero.}
\label{fig:diana}
\end{figure}

Para un calendario de quemas, que suma detecciones de muchos días y muchos píxeles, los errores
aleatorios pesan poco. Los sesgos lo son todo. Y la diferencia entre MODIS y VIIRS es,
precisamente, una diferencia de sesgos: ambos subestiman la actividad real, pero en distinta
medida.

\begin{parareto}
No hace falta eliminar los sesgos para armonizar; basta con que ambos registros tengan
\emph{el mismo}. La meta realista no es reconstruir <<el fuego verdadero>>, sino una serie cuyo
sesgo sea constante en el tiempo, de modo que sus variaciones sean reales. Hay que decirlo así
en la presentación: el calendario muestra actividad \emph{detectable}, no actividad total.
\end{parareto}

\section{Fuentes de omisión}

Reunimos aquí todo lo que hace que un fuego real no llegue al registro.

\begin{description}[leftmargin=1.2em,style=nextline]
  \item[Tamaño y temperatura.] El fuego no eleva lo suficiente la señal del píxel
  (capítulo~\ref{cap:deteccion}). Afecta mucho más a MODIS.
  \item[Nubes.] Una nube espesa oculta el fuego (capítulo~\ref{cap:atmosfera}). Afecta a ambos
  por igual en el mismo instante.
  \item[Hora.] El fuego no ardía cuando pasó el satélite (capítulo~\ref{cap:orbitas}). La
  duración importa: una quema de una hora tiene pocas posibilidades de coincidir con una pasada.
  \item[Geometría.] Píxeles grandes en el borde de la franja y huecos entre órbitas
  (capítulos~\ref{cap:orbitas} y~\ref{cap:escaneo}). Afecta sobre todo a MODIS.
  \item[Dosel y relieve.] Un fuego de sotobosque bajo copas densas, o en una ladera que da la
  espalda al sensor, queda parcialmente oculto.
  \item[Entorno.] Un fondo heterogéneo eleva los umbrales contextuales; un fuego rodeado de
  otros fuegos tiene pocos vecinos válidos.
\end{description}

\begin{ejemplo}{probabilidad de coincidir con una pasada}
Una quema agrícola dura en promedio 1,5 horas y se enciende en cualquier momento entre las 9:00
y las 17:00 con igual probabilidad. ¿Qué probabilidad hay de que esté ardiendo en una pasada a
las 13:30?
\solucion La quema está activa a las 13:30 si empezó entre las 12:00 y las 13:30: un intervalo
de 1,5~h dentro de las 8~h posibles. La probabilidad es $1{,}5/8 \approx 0{,}19$, un 19\,{}\%. Aun con un
sensor perfecto y cielo despejado, cuatro de cada cinco quemas de este tipo son invisibles para
una sola pasada diurna. Con tres satélites VIIRS que pasan dentro de un intervalo de unos 100
minutos, la ventana efectiva se alarga y la probabilidad sube aproximadamente a
$(1{,}5 + 1{,}7)/8 \approx 0{,}40$, un 40\,{}\%.
\end{ejemplo}

El ejemplo deja una lección de humildad sobre lo que mide un registro de focos, y explica por
qué añadir satélites a la misma órbita aumenta tanto las detecciones de quemas breves.

\section{Fuentes de comisión}

Menos frecuentes, pero concentradas en ciertos lugares, de modo que localmente pueden dominar.

\begin{description}[leftmargin=1.2em,style=nextline]
  \item[Brillo solar.] Agua, arrozales inundados, techos metálicos, invernaderos, paneles
  solares. Solo de día.
  \item[Suelos calientes y reflectantes.] Desiertos, salares, canteras, campos arados.
  \item[Bordes.] Costas, orillas, límites de nubes, claros recientes en el bosque.
  \item[Ruido del detector.] En VIIRS, sobre todo de noche y dentro de la Anomalía del
  Atlántico Sur.
\end{description}

\section{Fuentes térmicas persistentes}

Hay una categoría que no es ni omisión ni comisión, porque la detección es correcta: hay algo
muy caliente. Pero no es fuego de vegetación.

\begin{definicion}{fuente térmica estática}
Una fuente térmica estática es una anomalía térmica que aparece siempre en el mismo lugar por
causas distintas de la quema de vegetación: volcanes activos, antorchas de quema de gas en
pozos y refinerías, altos hornos, fundiciones, hornos de cemento, incendios de carbón
subterráneos.
\end{definicion}

Su firma es inconfundible en el tiempo: mientras que un fuego de vegetación dura horas o días y
no se repite en el mismo píxel hasta la temporada siguiente, una antorcha de gas genera
detecciones en la misma ubicación durante todo el año, de día y de noche
(figura~\ref{fig:estatica}).

\begin{figure}[ht]
\centering
\begin{tikzpicture}
\begin{axis}[width=13cm,height=5.4cm,ybar,bar width=9pt,ymin=0,ymax=32,
  symbolic x coords={E,F,M,A,My,J,Jl,Ag,S,O,N,D},xtick=data,
  xticklabels={E,F,M,A,M,J,J,A,S,O,N,D},
  ylabel={Días con detección},axis lines*=left,ymajorgrids,grid style={ceniza!20},
  legend style={at={(0.5,1.04)},anchor=south,font=\small,draw=none,legend columns=2,
    /tikz/every even column/.append style={column sep=10pt}}]
\addplot[fill=bosque!70,draw=none,area legend] coordinates {(E,0)(F,0)(M,0)(A,0)(My,1)(J,2)(Jl,4)(Ag,9)(S,7)(O,2)(N,0)(D,0)};
\addplot[fill=ceniza!80,draw=none,area legend] coordinates {(E,22)(F,19)(M,21)(A,24)(My,26)(J,27)(Jl,28)(Ag,27)(S,25)(O,23)(N,20)(D,21)};
\legend{Quemas de vegetación,Fuente estática}
\end{axis}
\end{tikzpicture}
\caption{Firma temporal de dos celdas (datos ilustrativos). La de vegetación tiene temporada; la
de la fuente estática registra detecciones casi todos los días despejados del año.}
\label{fig:estatica}
\end{figure}

Los datos archivados de FIRMS incluyen una columna \texttt{type} que clasifica cada detección
(capítulo~\ref{cap:firms}): 0 para presunto fuego de vegetación, 1 para volcán activo, 2 para
otra fuente estática en tierra y 3 para detecciones en el mar.

\begin{atencion}[Las fuentes estáticas contaminan el calendario]
Una sola antorcha de gas puede aportar más detecciones al año que todas las quemas de una
provincia, y las aporta de forma uniforme, aplanando la estacionalidad. Fíltralas con
\texttt{type = 0}. Como esa clasificación puede no estar disponible en los datos más recientes,
ten un segundo criterio propio: excluir las celdas con detecciones en, por ejemplo, más de 150
días distintos al año.
\end{atencion}

\section{Conteo múltiple}

Por último están los mecanismos que no añaden ni quitan fuegos, sino que hacen que un mismo
fuego aparezca varias veces. No son errores del algoritmo; son consecuencia de contar puntos.

\begin{table}[ht]
\centering
\renewcommand{\arraystretch}{1.2}
\begin{tabularx}{\textwidth}{>{\raggedright\arraybackslash}p{4.4cm}>{\raggedright\arraybackslash}Xcc}
\toprule
Mecanismo & Qué ocurre & MODIS & VIIRS \\
\midrule
Extensión espacial & un fuego grande ocupa varios píxeles & sí & sí, más \\
Solape entre píxeles vecinos & la huella real es más ancha que el espaciado & sí & menor \\
Efecto pajarita & barridos consecutivos se pisan lejos del nadir & sí & eliminado \\
Órbitas consecutivas & en latitudes altas las franjas se solapan & poco & sí \\
Varios satélites & cada satélite observa el mismo fuego & 2 & hasta 3 \\
Persistencia & el fuego sigue ardiendo en la pasada siguiente & sí & sí \\
\bottomrule
\end{tabularx}
\caption{Mecanismos por los que un mismo fuego genera más de un foco.}
\label{tab:multiple}
\end{table}

Una idea sencilla desactiva casi toda la tabla: no contar puntos, sino preguntar, para cada
celda de una cuadrícula y cada día, \emph{si hubo o no hubo detección}. Un fuego visto cinco
veces en una celda en un día cuenta lo mismo que uno visto una vez. Esta unidad, la
\termino{celda-día con fuego}, se desarrolla en el capítulo~\ref{cap:armonizacion}.

\section{Qué errores importan para armonizar}

No todos los errores afectan por igual a la comparación entre sensores. La
figura~\ref{fig:clasif-errores} los clasifica según dos preguntas: ¿afecta a ambos sensores por
igual? ¿cambia a lo largo del tiempo?

\begin{figure}[ht]
\centering
\begin{tikzpicture}[x=6.6cm,y=3.1cm]
  \fill[bosqueClaro] (0,1) rectangle (1,2); \fill[brasaClaro] (1,1) rectangle (2,2);
  \fill[brasaClaro] (0,0) rectangle (1,1); \fill[fuegoClaro] (1,0) rectangle (2,1);
  \draw[ceniza] (0,0) rectangle (2,2); \draw[ceniza] (1,0)--(1,2) (0,1)--(2,1);
  \node[font=\small\bfseries] at (0.5,2.15) {Igual en ambos sensores};
  \node[font=\small\bfseries] at (1.5,2.15) {Distinto en cada sensor};
  \node[font=\small\bfseries,rotate=90] at (-0.09,1.5) {Estable};
  \node[font=\small\bfseries,rotate=90] at (-0.09,0.5) {Cambia con el tiempo};
  \node[font=\small,align=center,text width=6cm] at (0.5,1.5) {Nubes en un instante dado\\ Fuegos bajo el dosel\\[3pt] \emph{No estorba: se cancela}};
  \node[font=\small,align=center,text width=6cm] at (1.5,1.5) {Tamaño del píxel\\ Crecimiento en los bordes\\ Huecos entre órbitas\\[3pt] \emph{Se corrige con calibración}};
  \node[font=\small,align=center,text width=6cm] at (0.5,0.5) {Nubosidad de cada año\\[3pt] \emph{Parte de la variabilidad real observable}};
  \node[font=\small,align=center,text width=6cm] at (1.5,0.5) {Número de satélites\\ Deriva orbital de Terra y Aqua\\ Versiones del algoritmo\\[3pt] \emph{Lo más peligroso: hay que controlarlo}};
\end{tikzpicture}
\caption{Clasificación de los errores según su efecto sobre la armonización.}
\label{fig:clasif-errores}
\end{figure}

El cuadrante inferior derecho es el que más cuidado exige. Una calibración entre MODIS y VIIRS
supone que la relación entre ambos es estable. Todo lo que la haga variar de un año a otro
(añadir un satélite, que una órbita derive, cambiar de versión de los datos a mitad de la serie)
rompe ese supuesto. De ahí las tres reglas prácticas que han ido apareciendo: fijar el conjunto
de satélites, calibrar con los años centrales del solape y usar una sola colección de datos para
todo el periodo.

\resumen

\begin{itemize}
  \item \textbf{Omisión}: fracción de fuegos reales no detectados. \textbf{Comisión}: fracción
  de detecciones que no son fuego. Los algoritmos son conservadores: poca comisión, mucha
  omisión.
  \item La omisión depende del tamaño del fuego; es una curva, no un número.
  \item Un \textbf{error aleatorio} se promedia; un \textbf{sesgo} no. Armonizar es igualar
  sesgos, no eliminarlos.
  \item Omiten fuegos: el tamaño, las nubes, la hora, la geometría, el dosel y el entorno.
  \item Las \textbf{fuentes estáticas} (volcanes, antorchas de gas, industria) son detecciones
  correctas que no son fuego de vegetación; se filtran con la columna \texttt{type} o por su
  persistencia.
  \item El \textbf{conteo múltiple} se neutraliza contando celdas-día con fuego en lugar de
  puntos.
  \item Los errores más peligrosos son los que difieren entre sensores \emph{y} cambian con el
  tiempo.
\end{itemize}

\ejercicios

\begin{enumerate}
  \item Un sensor reporta 420 detecciones en una zona donde había 600 fuegos; 12 de las
  detecciones eran falsas. Calcula omisión, comisión, probabilidad de detección y precisión.
  \item Un colega afirma que, como la comisión de VIIRS es del 1\,{}\%, <<VIIRS detecta el 99\,{}\%
  de los fuegos>>. Explica el error.
  \item Con las curvas de la figura~\ref{fig:pdet}, estima la probabilidad de detección de un
  fuego de \SI{100}{\metre\squared} y de uno de \SI{1000}{\metre\squared} para cada curva.
  \item Repite el ejemplo de la coincidencia con la pasada para un incendio forestal que dura
  30 horas. ¿Por qué el registro de focos representa mejor los incendios largos que las quemas
  breves?
  \item Para cada uno de estos casos, indica si es omisión, comisión, fuente estática o conteo
  múltiple: (a) un arrozal inundado aparece como foco al mediodía; (b) una refinería aparece 300
  días al año; (c) un incendio aparece con 40 focos en una pasada; (d) una quema bajo una nube
  no aparece; (e) un fuego aparece en Suomi NPP y en NOAA-20 con 50 minutos de diferencia.
  \item Coloca en los cuadrantes de la figura~\ref{fig:clasif-errores}: el brillo solar; la
  entrada en servicio de NOAA-21; una sequía que reduce la nubosidad.
  \item Diseña un criterio, usando solo fecha y coordenadas, para marcar como <<probable fuente
  estática>> una celda de 1~km. Discute qué fuegos reales podrías eliminar por error.
\end{enumerate}
